\begin{abstract}
Multimodal meme classifiers are evaluated on clean benchmarks, yet memes circulate in Romanised spellings, with mixed
languages and degraded images. We present ToonSense, a robustness study of frozen-encoder multimodal models for
Bangla--English memes, centred on cartoon memes from CMBan and replicated on BanglaAbuseMeme and MemoSen. We benchmark
six multilingual text encoders, four image encoders and five fusion heads with repeated cross-validation and
Holm-corrected tests, subject XLM-R + CLIP ViT-B/32 models to twelve linguistic and visual perturbations at three
severities, compare Integrated Gradients and Grad-CAM attributions before and after perturbation, and categorise
errors. Image features carry most of the signal, and fusion adds a small but significant gain on the larger datasets.
Romanised Bangla is the most damaging perturbation on every dataset, reducing fusion macro-F1 by 8--21\% relative,
whereas native-script spelling noise and code-mixing are nearly harmless. The perturbed text gains rather than loses
attribution share, and visual perturbations flip up to 27\% of predictions with small F1 changes. Training-time
augmentation with transliterated copies cuts the Romanisation drop to 2--5\%, including for an unseen spelling style,
at a clean-accuracy cost below one point.
\end{abstract}

\begin{IEEEkeywords}
Bangla memes, multimodal classification, code-mixing, robustness, Romanisation, CLIP, XLM-R, explainability
\end{IEEEkeywords}
